% ============================================================
%  SentinelAPI — Clean, non-technical explainer
%  Compile:  pdflatex explainer.tex
% ============================================================
\documentclass[11pt]{article}

\usepackage[margin=1in]{geometry}
\usepackage[T1]{fontenc}
\usepackage[utf8]{inputenc}
\usepackage{xcolor}
\usepackage{titlesec}
\usepackage{enumitem}
\usepackage{booktabs}
\usepackage{helvet}
\usepackage[hidelinks]{hyperref}

% clean sans-serif look
\renewcommand{\familydefault}{\sfdefault}

% accent colour
\definecolor{accent}{RGB}{37, 99, 235}   % blue
\definecolor{red}{RGB}{220, 38, 38}
\definecolor{orange}{RGB}{234, 88, 12}
\definecolor{grayline}{RGB}{229, 231, 235}

\titleformat{\section}
  {\Large\bfseries\color{accent}}{\thesection.}{0.6em}{}
\titleformat{\subsection}
  {\large\bfseries}{\thesubsection}{0.6em}{}

% simple "card" box
\newcommand{\card}[2]{%
  \noindent\fcolorbox{grayline}{white}{%
    \begin{minipage}{\dimexpr\linewidth-2\fboxsep-2\fboxrule\relax}
      \textbf{\color{accent}#1}\\[2mm]
      #2
    \end{minipage}}%
  \par\vspace{3mm}}

\title{\bfseries SentinelAPI\\[1mm] \Large What we built, how it works, and why it matters}
\author{AMIHACKS 2026 --- Track C (Deep-Tech)}
\date{\today}

\begin{document}

\maketitle
\thispagestyle{empty}

\vspace{2mm}
\begin{center}
\fcolorbox{accent}{white}{\parbox{0.92\linewidth}{%
\centering\itshape
``We built a tool that finds a website's hidden security holes automatically ---
before a hacker does.''}}\end{center}

\vspace{6mm}

\section{The Problem}

Apps talk to servers through ``doors'' called \textbf{APIs}. Every time you log in,
check a balance, or open a profile, your app is knocking on one of these doors.
When a door is built carelessly, two things go wrong:

\begin{enumerate}[leftmargin=1.6em]
  \item \textbf{The open door.} The server hands over \emph{someone else's} data
        when you ask --- because nobody checked who is asking.
  \item \textbf{The oversharing door.} The server sends back \emph{too much} ---
        passwords, card numbers, private details the app never needed.
\end{enumerate}

Hackers find these mistakes every day. Companies usually notice only \emph{after}
the damage is done, because checking by hand is slow and expensive.

\section{What We Built}

\textbf{SentinelAPI} is an automatic security inspector. You point it at an app's
doors, and it tells you exactly which ones are unsafe --- in plain words, with a
one-line proof, so anyone can act on it.

\card{In one sentence}{
It creates two imaginary users, checks whether one can see the other's secrets,
and checks whether the server is sharing more than it promised.}

\section{How It Works}

\noindent Five simple steps:

\begin{enumerate}[leftmargin=1.8em, itemsep=2mm]
  \item \textbf{We built a practice app} with known flaws, like a training dummy
        for the inspector to test itself on.
  \item \textbf{The inspector creates two imaginary users} --- let's call them
        Alice and Bob.
  \item \textbf{It pretends to be Alice and asks for Bob's stuff.}
        If the server says \emph{``sure, here's Bob's data''} --- that's a red flag.
        We call it the \emph{open-door} problem.
  \item \textbf{It compares the promise with reality.} The server's own rulebook
        says what it will share; the inspector checks what it \emph{actually} shares.
        Extra secrets --- passwords, card numbers --- are another red flag. We call
        it the \emph{oversharing} problem.
  \item \textbf{It writes a report.} Every problem gets a colour badge
        (\textcolor{red}{red} = serious, \textcolor{orange}{orange} = important)
        plus a copy-paste command that proves it.
\end{enumerate}

\section{The Tech Stack}

\noindent The building blocks, in plain terms:

\begin{center}
\begin{tabular}{ll}
\toprule
\textbf{Tool} & \textbf{What it does for us} \\
\midrule
Python + FastAPI & The practice app, and the ``brain'' of the inspector \\
httpx & How the inspector knocks on the app's doors \\
OpenAPI & The server's rulebook the inspector reads \\
HTML + CSS & The clean report page with the colour badges \\
Rust (planned) & A faster inspector engine for very large apps \\
\bottomrule
\end{tabular}
\end{center}

\section{What's New About It}

\card{We test, we don't guess}{
Most scanners throw random junk at a server and hope something breaks. We do
something smarter: we create \textbf{two real users} and check whether one can see
the other's data. That is a \emph{real test}, not a guess.}

\card{We use the server's own rulebook}{
We compare what the server \emph{promised} to share against what it actually
shares. The difference between promise and reality is where the leak is.}

\card{Every finding comes with proof}{
No vague warnings. Each problem includes a copy-paste command anyone can run to
see the leak with their own eyes --- so there are \emph{no false alarms}.}

\card{Runs anywhere, costs nothing}{
The whole thing works offline with free, open tools --- no paid services, no
cloud, no internet needed for the demo.}

\card{Built to grow}{
The engine is designed so we can swap the Python brain for a faster \textbf{Rust}
one later, without changing anything else.}

\section{How We Show It}

\noindent The live demo takes about \textbf{30 seconds}:

\begin{enumerate}[leftmargin=1.8em, itemsep=2mm]
  \item Start the practice app.
  \item Run the inspector.
  \item Open the report --- a clean page with \textcolor{red}{red} and
        \textcolor{orange}{orange} cards, each one saying \emph{what} is wrong,
        \emph{why}, and \emph{how to prove it}.
  \item Run the proof command on screen --- and watch Bob's private data appear
        while logged in as Alice.
\end{enumerate}

\vspace{4mm}
\begin{center}
\fcolorbox{grayline}{white}{\parbox{0.92\linewidth}{%
\centering\large\bfseries
The result: a tool that finds security holes automatically, explains them in
plain words, and proves them --- before a hacker does.}}\end{center}

\end{document}
